\section{Limitations and Future Work}
\label{sec:limitations}

The preceding sections establish the Bailout Protocol as a structurally compulsive escalation architecture with formally stated Safety, Liveness, and Accountability properties. These guarantees hold under a set of explicit platform and deployment assumptions; they degrade or dissolve outside those assumptions. This section makes those assumptions concrete as four honest limitations, each paired with a directed research agenda that maps directly to the three future-work items that follow. Together, limitations and future work define the boundary between what the protocol achieves and what subsequent research must establish.

\subsection{Limitations}

\subsubsection{Human Response Latency.}
\label{sec:limitations-latency}

The Liveness property (Theorem~\ref{thrm:termination}) guarantees eventual human contact but imposes no bound on the time horizon in which that contact is operationally useful. The bound $T_{\max}$ derived in the liveness proof is a function of network topology and retry parameters, not of human cognitive response time. In high-frequency domains---algorithmic trading, real-time process control, autonomous vehicle navigation---the time required for a human principal to receive, interpret, and respond to an escalation may exceed the operational time constant of the system by an order of magnitude.

The protocol's failure mode under latency is well-defined: the system stalls in \texttt{ESCALATING} until the human responds. This is the intended behavior, but the operational cost of the stall is domain-dependent. In a precision agriculture context such as Irrig8, a stall delaying a watering decision by minutes or hours is acceptable. In a financial risk management context, a stall delaying a risk-interlock response by even seconds may be economically catastrophic. The protocol currently provides no mechanism for expressing domain-specific cost structure or adapting stall-tolerance accordingly. Timing parameters are provisioned statically at node setup and are not sensitive to the anchor's operational context: a human supervising a single node in normal conditions may respond within seconds; the same human supervising a fleet during an emergency may have queues taking minutes to clear. The protocol has no mechanism to distinguish between these scenarios or to trigger a faster escalation path when a secondary anchor is available.

\textbf{Directed research:} Empirical measurement of human response latency distributions across deployment contexts to establish statistically grounded latency budgets per domain. Integration of these budgets as enforceable policy parameters at node provisioning, with the Bounds Engine calibrated to the provisioned budget. Investigation of graded escalation responses: a protocol that requests a rapid acknowledgment (within seconds) followed by a detailed response (within hours) better matches the operational reality of human decision-making under varying cognitive loads.

\subsubsection{Adversarial Conditions.}
\label{sec:limitations-adversarial}

The Safety property (Theorem~\ref{thrm:drift-prevention}) requires that human-contact pathways cannot be closed without recorded human authorization. This guarantee holds under the assumption that the host platform provides independent scheduling priority to the Bailout Monitor (Assumption A1 in Section~\ref{sec:properties}). In an adversarial environment---particularly one in which the host platform itself may be compromised---this assumption may not hold. A compromised OS or virtualized substrate can modify thread priorities, preempt the monitor thread, or introduce race conditions into the pathway-closure operations the monitor is designed to interdict.

The protocol's guarantees are platform guarantees. They are derived from invariants the host platform must provide: non-preemptible monitor scheduling, append-only storage for the authorization ledger, and interruptible blocking calls for pathway management. If the platform itself is adversarial, these invariants cannot be relied upon. A sophisticated attacker who controls the execution substrate can rewrite the Bailout Monitor's code, suppress its scheduling, or tamper with the authorization ledger before persistence. The bypass mechanism enforces the no-machine-actor property on the propagation path \emph{once escalation is triggered}, but it cannot protect against a machine actor who controls the substrate on which the propagation path runs.

The applicability boundary is deployments where the host platform is trusted but individual agents may behave unexpectedly: the standard assumption for production AI in regulated industries. It is not designed for deployments in which the platform itself is an active adversary. For national infrastructure, military systems, or financial clearinghouses, the protocol must be composed with substrate-integrity mechanisms outside its current scope: secure boot chains, hardware-attested execution environments, and independent watchdog monitors on separate hardware.

\textbf{Directed research:} Empirical evaluation of Bailout Monitor scheduling guarantees under mainstream hypervisor and OS kernel configurations. Specification of a minimal substrate-integrity interface that the protocol requires from the host platform, enabling formal composition arguments with orthogonal integrity mechanisms.

\subsubsection{Scalability Under 1:1 Anchoring.}
\label{sec:limitations-scalability}

The termination guarantee (Theorem~\ref{thrm:termination}) bounds hierarchy \emph{depth} but not \emph{breadth}. Under the current 1:1 anchoring model, a single human anchor $p \in P$ may supervise many independent subtrees, each spawning children to depth $d_{\max}$. As active node count grows, spawn events, ledger entries, and inter-node communication grow proportionally. The forensic ledger's append-only property implies monotonic growth without compaction: no garbage collection or archival policy is specified.

For deployments at the scale envisioned for Irrig8 across large agricultural regions---thousands of edge nodes across tens of thousands of acres---this growth is a practical operational concern. Cross-hierarchy coordination (Corollary~\ref{corol:compositional-safety}) requires explicit Purpose Layer authorization recorded in each participating hierarchy's ledger; as independent hierarchies multiply, coordination overhead at a single anchor creates a bottleneck the protocol does not address. The 1:1 anchoring model is a structural proof-of-concept. Enterprise and infrastructure deployment requires at minimum 1:N pooled supervision, and potentially N:M anchor pools with load balancing, before the protocol can be certified at scale.

\textbf{Directed research:} Benchmarks on forensic ledger growth rates under simulated production workloads at 10$\times$, 100$\times$, and 1000$\times$ current Irrig8 deployment scale. Formal specification of a ledger compaction and archival policy that preserves audit integrity while enabling operational sustainability. Design of cross-hierarchy coordination primitives that reduce anchor coordination overhead without relaxing Safety.

\subsubsection{Exploitation Risk from Misconfigured Bounds.}
\label{sec:limitations-exploitation}

The Bounds Engine enforces depth limits and scope predicates as configured parameters. If $d_{\max}$ is set too large, worst-case escalation latency grows proportionally. If scope predicates are miscalibrated---too permissive---the bailout triggers so rarely that the protocol becomes a formality; too conservative, and alert fatigue produces operator habituation, causing genuine bailout events to be misclassified as noise.

More insidiously: a sophisticated agent that understands threshold calibration may strategically set scope boundaries just inside the bailout trigger, absorbing exception conditions that should escalate without actually triggering escalation. This is not a code defect; it is a structural vulnerability in any bounds-based enforcement model. The protocol detects escalation \emph{once it has been triggered}, but it does not detect whether a machine actor has intentionally avoided triggering in the first place. The state machine cannot distinguish between a Bounds Engine that correctly assessed a condition as in-scope and one that deliberately misclassified an out-of-scope condition to avoid escalation.

\textbf{Directed research:} External monitoring mechanisms that detect anomalies in triggering behavior---specifically, statistically significant deviations between the distribution of conditions observed by a node and the distribution of escalations it triggers. Entropy-based trigger monitors: a node that observes conditions outside its operating range but triggers escalation at a rate statistically inconsistent with its observed condition distribution is a candidate for audit. Investigation of provenance attestation for operating-range assessments: can the protocol require that each Bounds Engine assessment be counter-signed by an independent audit process recording the raw inputs and assessment criteria, making post-hoc misclassification detectable?

\subsection{Future Work}

\subsubsection{Formal Verification.}
\label{sec:future-formal-verification}

The correctness properties in Section~\ref{sec:properties} are supported by proof sketches using argument maps. Full formal verification using a proof assistant (Coq, Lean, or Isabelle/HOL) or a formal specification language (TLA$^+$) would close gaps in the informal arguments, confirm the logical consistency of the three-layer architecture, and produce a machine-checkable artifact suitable for regulatory review.

Of particular interest is whether the functional separation between Purpose, Bounds, and Fact layers admits a compositional verification argument: if each layer satisfies its specification, does the composition necessarily satisfy the protocol's guarantees? A compositional verification architecture would distribute the verification burden across layer implementations rather than centralizing it on the full system, and would make the protocol more robust to implementation changes. Equally valuable: formal modeling of the immutable parent pointer $\alpha$ as a hardware-protected function, and a formal proof that the forensic ledger's append-only property is preserved under concurrent writes from multiple Fact Layer processes in a distributed deployment.

Formal verification would expose hidden assumptions in the current argument-map proofs---gaps that appear correct but lack rigorous justification. The effort would thus be as valuable for what it reveals about the protocol's limitations as for what it confirms about its correctness.

\subsubsection{Adaptive Timeout.}
\label{sec:future-adaptive-timeout}

The fixed timeout model is a conservative baseline that optimizes for worst-case deployment rather than typical-case efficiency. Adaptive timeout mechanisms should adjust the escalation window based on: current depth (nodes near $d_{\max}$ warrant higher urgency than nodes at depth 1); node criticality classification (physical actuator nodes such as irrigation valves warrant faster escalation than informational nodes); historical anchor response patterns (a consistently responsive anchor may be assigned shorter windows, with automatic fallback to a secondary anchor); and environmental conditions (hazardous readings may trigger preemptive escalation before timeout, trading false-positive cost against catastrophic failure cost).

Adaptation requires a feedback loop updating timeout parameters from observed operator response data---a control theory problem amenable to a proportional-integral-derivative (PID) controller or a lightweight reinforcement learning agent. The critical constraint: adaptation must preserve the deterministic timing guarantee required by the Safety property. The applicable timeout for any given bailout signal must be computable from publicly observable state, preventing the timeout value itself from becoming an attack surface. The adaptive timeout must also be recorded in the forensic ledger, ensuring that changes to escalation timing remain auditable after the fact.

\subsubsection{Distributed Human Anchor Pools.}
\label{sec:future-anchor-pools}

The current specification assumes a fixed 1:1 pairing between spawned node and human anchor. Enterprise and infrastructure deployments require a pool model: a set of human anchors $P'$ who share responsibility for a set of nodes, with a routing policy that delivers escalations to available anchors within $T_{\max}$ under high concurrent load. This implies load-aware routing rather than static round-robin. The pool architecture must preserve the Accountability property: every escalation must be attributed to a specific human anchor, immutably recorded, requiring that routing decisions are committed to the forensic ledger. The pool architecture must also include a load-shedding policy---a defined behavior for when all anchors in a pool are simultaneously overloaded---that preserves Safety while providing meaningful feedback about the overload condition.

Open questions include: cross-jurisdictional authorization (an anchor in one jurisdiction may lack legal authority to adjudicate actions in another, requiring the accountability chain to record which anchor resolved an escalation and the jurisdictional scope of their authorization); collusion risk among pool members authorizing depth budgets that collectively violate safety properties; and the treatment of institutional actors---corporate compliance teams or designated officers---as anchor pool members, requiring institutional rather than individual operational protocols.

\section{Conclusion}
\label{sec:conclusion}

The Bailout Protocol advances a single, architecturally precise claim: that autonomous systems built on the \bxthree{} Framework must propagate every unresolved exception state to a named human principal, bypassing all intermediate machine actors, as a compulsory structural requirement rather than a runtime choice. This mandatory bailout property is not a feature that can be retrofitted to an existing agentic architecture through better error handling or improved escalation heuristics. It is a structural property that follows, as a matter of architectural necessity, from the functional separation between the Purpose, Bounds, and Fact layers. When the Bounds Layer encounters a condition outside its provisioned operating range that it cannot resolve, no machine actor holds the authority to adjudicate that condition. The decision must return to the Purpose Layer, which carries intent, judgment, and final authority. The Bailout Protocol is the mechanism that enforces this return---compulsorily, deterministically, and with full forensic attribution at every step.

This contribution is realized within the \bxthree{} Framework as its fifth named architectural pillar, complementing Loop Isolation, Recursive Spawning, Spatial Firewall, and Sandbox Gate. Together, these five pillars constitute the upstream accountability guarantee: that \bxthree{} systems \emph{fail upward into human consciousness, never downward into autonomous chaos}. The Bailout Protocol is the runtime expression of this guarantee for all exception conditions not resolved by the preceding four pillars---including conditions not anticipated at design time, which is precisely the operational reality that makes the protocol necessary.

The \bxthree{} Framework's role in this contribution is to provide the conceptual architecture within which the Bailout Protocol operates as a necessary consequence rather than an independent design choice. The functional separation of Purpose, Bounds, and Fact layers creates the structural conditions that make the bailout requirement compulsory. The Forensic Ledger maintains the immutable attribution record that makes every escalation auditable. The immutable parent-pointer established at node provisioning ensures that the escalation path is fixed at the architectural level and cannot be modified by any machine actor. These elements are not independent mechanisms that happen to cooperate; they are mutually defining structural facts whose necessity is established by the protocol's own specification. The upstream accountability guarantee and the Bailout Protocol are two aspects of the same architectural commitment: that in a \bxthree{} system, every exception that exceeds machine authority reaches a human who can bind it, and every such exception is recorded so that the binding can be verified.

The theoretical claim is foundational: that the upstream accountability guarantee of the \bxthree{} Framework is not a design aspiration but a logical consequence of the framework's architecture, and that the Bailout Protocol is the mechanism by which that consequence is enforced at runtime. The practical claim is equally clear: that autonomous systems operating without such a mechanism are structurally exposed to failure modes---invisible exception absorption, contingent escalation, and irreversible drift---that cannot be eliminated by incremental improvement to existing error-handling or oversight practices. They require a new architectural primitive. The Bailout Protocol is that primitive.

The limitations identified in Section~\ref{sec:limitations} define the boundary between what the protocol achieves within its scope and what must be established by subsequent work. Human response latency, adversarial conditions, scalability, and exploitation risk are constraints on the deployment context, not flaws in the protocol's logical design. They are the research agenda for the work that follows: formal verification to establish the correctness proofs rigorously; adaptive timeout to reduce provisioning burden and improve operational fit across deployment contexts; and distributed human anchor pools to enable deployment at industrial scale while preserving the Safety, Liveness, and Accountability guarantees. The protocol's structural necessity within the \bxthree{} Framework is established by its specification. Its practical necessity in the field is established by the operational reality of autonomous systems that, without it, have no architectural path to human accountability when they encounter conditions that exceed machine authority.
